\needspace{5cm} \section{Manuál pomocného rozhodčího} \label{manuál pomocného rozhodčího} Role \odkaz{pomocného rozhodčího}{pomocný rozhodčí} bývá při \odkaz{zápasech}{zápas} podceňována, a bývají do ní obsazováni nejméně zkušení \odkaz{hráči}{hráč}. Bývají dva pomocní rozhodčí. 
 
Obvykle jsou pomocňáci třeba již před představením  k~přípravě sálu  - rozmisťování židliček, přípravě \odkaz{hlasovacích kartiček}{hlasovací kartička} a \odkaz{papučí}{papuče} na jednotlivá místa, 
rozvěšením bodovacích tabulek, pomocí \odkaz{technikovi}{osvětlovač} s~ozkoušením mikrofonů apod. 
 
Často také prodávají lístky - pak je nutné diváky přivítat, prvodivákům vysvětlit lístečky s~tématy a předat hlasovací karty a papuče (pokud již nejsou připravené na jednotlivých místech). 
 
Těsně před zápasem pak od diváků vybírají napsaná témata a dávají je do košíku na témata. 
 
Na scénu přichází na vyzvání konferenciéra těsně před \odkaz{hlavním rozhodčím}{rozhodčí} - krátce poté, co se hráči vrátili z~\odkaz{veřejné rozcvičky}{veřejná rozcvička}. 
 
Jejich hlavní rolí je pomoci rozhodčímu a hráčům k~hladkému průběhu zápasu. Proto je vhodné, když se ještě před zápasem domluví: 
\begin{itemize}
\item  kdo z~nich bude \odkaz{ukazovat čas}{ukazování času}  (ten by se to měl rychle doučit)
\item  kdo bude nosit košíček s~tématy
\item  a kdo bude kterému \odkaz{týmu}{tým} počítat hlasy diváků, a následně přidávat \odkaz{body}{bod} či \odkaz{trestné body}{trestný bod} za \odkaz{fauly}{faul}.
\end{itemize}
 
Po celou dobu zápasu pomocňák setrváva v~charakteru - v~postavě, která nadevše podporuje hlavního rozhodčího (HR) - takže tedy vrhá se do cesty papučí vrhaných na HR, mračí se na hráče, kteří s~HR nesouhlasí apod. To mu ovšem nebrání hrát drobné \odkaz{statusové}{status} hry s~druhým pomocňákem, a v~reakcích na vše, co se děje, si budovat vlastní osobnost. Pokud ovšem nejsou zapotřebí, měly by se odebrat někam stranou, a nebrat \odkaz{fokus}{fokus} ostatním na scéně. 
 
Pokud HR určí čas na kategorii ( nebo u~kategorií, kde je čas vždy např. - \odkaz{Smrt v~jedné minutě}{smrt v~jedné minutě},\odkaz{Polovina času}{polovina času}, \odkaz{Prskavka}{prskavka},  \odkaz{Ryba (ve čtyřech)}{ryba (ve čtyřech)}, \odkaz{Stíhačka}{stíhačka}, \odkaz{Dvojitá stíhačka}{dvojitá stíhačka}, \odkaz{Trojitá stíhačka}{trojitá stíhačka} (u~stíhaček je ukazováni času poněkud speciální, doporučujeme prostudovat)), jeden pomocňák ukazuje čas, ovšem toto ukazování je informací pro hráče a HR, teprve hvizd píšťalky HR kategorii ukončuje. 
 
Ve chvíli kdy kategorie skončila, HR obvykle přiděluje trestné body za fauly, pomocňáci je promptně přidávají na skórovací tabuli a v~případě, že jeden tým nasbírá tři trestné body, přidělí místo nich  druhému týmu jeden kladný bod, a informují hlavního rozhodčího, který vše může vysvětlit divákům. 
 
Následuje divácké hlasováni, kdy každý pomocňák sčítá jednu barvu a pak oznamuje výsledek HR, ten pak divákům (i pomocňákům) oznamuje, který tým získal kladný bod. Pomocňáci promptně bod týmu přidělí. 
 
Ve chvíli, kdy HR požádá o~hlasovací témata (často žádá jen nataženou paží), pomocňák pověřený nošením košíčku měl by témata ihned HR přinésti. 
\needspace{5cm} \section{Manuál konferenciéra} \label{manuál konferenciéra} Nastala ta velká chvíle, kdy budeš třeba poprvé konferovat \odkaz{Zápas}{zápas}. Tento manuál by ti měl pomoci se na to trochu připravit. 
 
\subsection{ Co dělat týdny předem } Nauč se přehledně popsat a vysvětlit všechny \odkaz{zápasové kategorie}{:kategorie:zápasové kategorie}. Co je důležité, co potřebují diváci vědět?  
\subsection{ Co dělat dny předem } Připrav si postavu. Rozhodni se, zda budeš konferovat jako 
\begin{itemize}
\item  Tematická postava (princezna, vodník, zedník)
\item  Alterego v~nějaké  specifické emoci/nastavení (zmatený, zamilovaný, ...)
\item  Čistě profesionální improvizátor.
\end{itemize}
 
Vyzkoušej si popsání zápasových kategorií v~rámci postavy. Jaká slova, jaké metafory by taková postava použila?  
 
Najdi či sežeň kostýmní doplňky, boty, rekvizity apod. 
 
\subsection{ Co dělat těsně předem } Ujasni si jména všech \odkaz{hráčů}{hráč}, \odkaz{pomocňáků}{pomocný rozhodčí}, \odkaz{osvětlovačů}{osvětlovač}, \odkaz{hudebníků}{hudebník} apod. Ideálně si je napiš, ať se nedopustíš faux pas. 
 
Vyber s~hráči \odkaz{vhodné kategorie}{:kategorie:rozcvičky} na veřejnou rozcvičku, a ujasni si, co na ně budeš potřebovat od diváků. 
 
Prober s~hudebníkem a osvětlovačem tvůj příchod na scénu. Má zkusit zahrát Imperial March, Vjezd gladiátorů, svatební pochod Mendelssohn, nebo Holku modrookou? 
 
Domluv s~garantem akce a rozhodčím, jak dlouhý má zápas být, kdy se má optimálně končit. 
 
\subsection{Zápas začíná} Přicházíš na scénu, vítáš se z~diváky a možná trošku představuješ svoji postavu. 
Nejspíš diváky trochu \odkaz{rozehřeješ}{rozehřívání diváků} a zjistíš, kdo je na zápase poprvé. 
 
Je vhodné divákům teď nebo po veřejné rozcvičce sdělit pár faktů/praktických  informací, např: 
\begin{itemize}
\item  co jsou zápasy
\item  jak fungují \odkaz{hlasovací kartičky}{hlasování}
\item  co je kazoo, píšťalka, \odkaz{faul}{faul}
\item  používání  míčků či papučí
\end{itemize}
 
Pokud něco z~toho opomeneš, nevadí, hlavní rozhodčí si to stejně nejspíš zkontroluje. 
 
\subsubsection{ Veřejná rozcvička } Zveš hráče na  \odkaz{veřejnou rozcvičku}{veřejná rozcvička}, hráči jsou ještě v~nediferencovaných černých trikotech, zatím je nemusíš představovat.  
Je dobré mít ujasněno, jaké \odkaz{Askfor}{askfor} pro rozcvičku potřebuješ, (např. prostředí pro \odkaz{Souboj otázek}{souboj otázek}) nebo jaký je průběh (třeba u~\odkaz{Stroje}{stroj} či \odkaz{Běhu přes překážky}{běh přes překážky} ) 
Sleduj čas, pocitově deset minut rozcvičky už je asi moc. 
 
Posíláš hráče se převléci do týmových dresů (čili aspoň chvíli na to potřebují), takže zatím můžeš vysvětlit další část z~informačního bloku. 
 
A~pak už zveš hráče na scénu, jmenovitě nejprve jeden tým jménem, a jmény jeho hráčů, pak druhý. 
Následně zveš pomocné rozhodčí a pak s~malou dramatickou paouzou hlavního rozhodčího a předáváš mu hlavní slovo. 
  
Od této chvíle sloužíš jako podpůrný prvek hlavního rozhodčího. Na jeho vyzvání vysvětluješ kategorie apod, a v~případě že nejsi na scéně zapotřebí, mizíš. 
Na vyzvání vyhlašuješ přestávku (na deset minut, takže za patnáct minut se tu všichni sejdeme apod). 
 
Druhou polovinu zahajuješ velmi krátce, jedna věta, a zveme hráče, pomocňáky, rozhodčího. 
Na úplný závěr, po výzvě k~ukončení zápasu si bereš mikrofon (ať jsi slyšet přes potlesk) a oznamuješ : 
 
\begin{itemize}
\item  Dnešní zápas vyhrál tým XY nebo skončil remízou za stavu x:y.
\item  Za tým X hráli jmenovitě ..
\item  Za tým Y hráli jmenovitě ..
\item  pomocňákovali, rozhodcovali,svítili, hudebnili, střihali video na dabing, tvořili grafiku plakátů apod ...
\item  pozvánka na další zápasy ať už na stránkách týmů nebo na www.improliga.cz/kalendar-predstaveni/
\item  klanění, odchod do šaten
\end{itemize}
\needspace{5cm} \section{Manuál rozhodčího} \label{manuál rozhodčího} Už nastala ta velká chvíle, kdy budeš pískat \odkaz{Zápas}{zápas} 
 
\subsection{ Kdo jsi } \begin{itemize}
\item  Jsi zkušený \odkaz{improvizátor}{hráč}, který zažil řadu zápasů
\item  Práce s~publikem  jako \odkaz{MC}{mc} či \odkaz{konferenciér}{konferenciér} ti není cizí.
\item  Rozpoznáváš \odkaz{fauly}{:kategorie:fauly} či jiné hřích proti improvizaci a umíš je dát najevo improvizátorům pomocí domluvených gest i verbálně
\end{itemize}
 
\subsection{ Co děláš před představením } Pravděpodobně vedeš \odkaz{Předzápasový trénink}{předzápasový trénink}. 
 
Zapamatuj si a zapiš jména hráčů, dohodni s~nimi improvizační kategorie, které ovládají. 
 
S~Garantem akce si ujasni plánovanou délku zápasu, nutný konec apod, případně techniku (např. na \odkaz{Vysílání FM}{vysílání fm} nebo \odkaz{Dabing filmu}{dabing filmu} 
 
\subsection{ Zápas začíná } Po nějaké době jsi konferenciérem pozván na scénu, zdravíš se s~ním a s~diváky. 
Pravděpodobně dovysvětluješ co konferencér nevysvětlil či to alespoň kontroluješ. 
Necháváš hráče zakřičet své pokřiky, případně slib kapitánů. 
A~pak už necháváš hráče hrát kategorie, ideálně vybrané tak, aby ty podobné nebyly vedle sebe, aby byl zápas dynamičtější. 
 
\subsection{ Průběh kategorie } Oznámíš kategorii, vybíráš téma, či jiné \odkaz{Askfor}{askfor} a oznamuješ čas. 
Pískáš začátek, a jsi připraven pískáním faulů, či zastavením improvizace a režijní editací, či komunikací s~hráči improvizaci vylepšovat. 
Pískáš konec kategorie, oznamuješ písklé fauly a přidělené \odkaz{trestné body}{trestný bod}. 
Vyhlašuješ hlasování a na základě informace od \odkaz{pomocných rozhodčí}{pomocný rozhodčí} přiděluješ \odkaz{body}{bod}. 
 
\subsection{ Přestávka } Ve vhodnou chvíli vyhlašuješ přestávku, v~zákulisí kontroluješ a vyhlašuješ návrat na zápas, kde se zase hrají kategorie. 
 
\subsection{ Druhá půlka } Hrané kategorie směřuješ k~nějaké typicky energeické pro všechny -  častá bývá \odkaz{Zpívaná}{zpívaná} nebo třeba \odkaz{Filmové žánry}{filmové žánry} zakončené písní. 
 
Vyhlašuješ výsledné skóre, konferenciér zve na scénu vítěze, klanečka všech účastníků, a odchod s~nimi do zákulisí. 
 
Tam s~velkou pravděpodobností vedeš pozápasovou reflexi, někde se ale nejprve uklízí prostory/loučí se s~diváky. 
 
\subsection{ P.S. } \begin{itemize}
\item  Nezapomeň, jsi vládcem scény.
\item  \odkaz{pomocní rozhodčí}{pomocný rozhodčí} jsou tu pro tebe.
\end{itemize}
